\section{Introduction}

Cardinality estimation is a small interface with system-wide consequences. A
query optimizer uses estimated result sizes to compare access paths, join
orders, and operator implementations; an estimate that is wrong by orders of
magnitude can therefore induce a bad plan even when the execution engine is
otherwise sound \cite{Leis2015HowGood}. PostgreSQL's ordinary statistics are
deliberately compact and are usually maintained per column. They consequently
cannot represent every cross-column correlation that appears in a workload.

Learned cardinality estimators have shown that this accuracy gap can be
substantial. The Are We Ready study also makes clear that accuracy is only one
part of the deployment question: training or update cost, inference cost,
staleness, interpretability, and behavior under changing data all matter
\cite{Wang2021AreWeReady}. This paper does not argue that learned estimators
are unsuitable for production. Instead, it asks whether a less disruptive
intervention can recover useful accuracy before the estimator itself is
replaced.

\emph{How far can native statistics take cardinality estimation?} More
precisely, instead of replacing PostgreSQL's cardinality estimator, we
investigate how much accuracy can be recovered by automatically designing the
native extended statistics that the deployed estimator already understands.
PostgreSQL supports multivariate statistics for correlations and functional
dependencies, but deciding which column groups to create, how to prioritize
them, and how to validate their benefit remains a physical and statistical
design problem \cite{PostgreSQL16Docs}.

\advisor approaches this problem as an offline, planner-in-the-loop workflow.
It freezes a portable schema, a fixed typed sample, and a workload snapshot;
constructs native statistics on the sample; evaluates catalogless hypothetical
native statistics in an isolated patched PostgreSQL planner; searches under a
bounded budget; and emits a DBA-reviewable recommendation that can be deployed
using stock PostgreSQL \texttt{CREATE STATISTICS} and \texttt{ANALYZE}. The
patched server is a design-time what-if substrate, not a required production
runtime. We use the phrase ``patched at design time, stock at deployment time''
as a working description of this separation, not as a claim that the workflow
is already production-ready.

The research focus is not a new cardinality estimator, a first statistics
advisor, or a novel greedy algorithm. The contribution is the combination of
five testable system and empirical questions:

\begin{itemize}
  \item a catalogless what-if substrate for PostgreSQL native MCV and
  functional-dependency payloads;
  \item a production-oriented architecture that isolates exploration and
  preserves a stock, DBA-controlled deployment path;
  \item a sample-to-full-data transfer protocol that separates design on a
  fixed sample from independent native-statistics recomputation on full data;
  \item an empirical study on the AreCELearnedYet benchmark lineage that
  revisits how much of the reported learned-CE gap native statistics can
  address; and
  \item an explicit evaluation of fidelity, ablations, and operational cost,
  rather than accuracy alone.
\end{itemize}

The central claims are intentionally conditional. A recommendation optimizes
the evaluated membership, not necessarily the union of that membership with
arbitrary externally managed production statistics. Existing production
extended statistics remain outside advisor ownership: the advisor does not
silently drop, rename, replace, or reconcile them. A DBA decides whether such
objects remain, are removed, are replaced, or coexist. Likewise, current
Census13, Forest10, Power7, and DMV11 evidence is pilot or preliminary until
the corresponding immutable artifacts are regenerated and audited.

The rest of the paper defines the design problem, describes the architecture
and PostgreSQL extension, gives an evaluation protocol for five research
questions, and records the result schema without manufacturing final numbers.
